\section{Extremal weights and a dominated base}\label{har:sec:extremal}

This route uses a source component for a one-parameter action rather
than the invariant index. It applies to a smooth projective rationally
connected \(S\) with smoothly semipositive \(L=-K_S\). The final
product argument will give invariance under an entire algebraic torus.
The source-character theorem of Section~\ref{inj:extremal-weights}
already gives this conclusion under weaker hypotheses. Here the
purpose is to obtain the character by harmonic vanishing: we construct
a base dominated by \(L\), then retract that base to its source so
that Euler interpolation takes place in one specified weight.

\begin{lemma}[Source weights and retraction]\label{har:lem:source}
Let \(S\) be smooth connected projective with a \(\mathbb C^*\)-action,
allowing the trivial action. Let \(S_0\) be its attracting component
for \(t\to0\), and write \(w_C(U)\) for the fiber weight of a
linearized line \(U\) on a fixed component \(C\).
A nonzero eigensection of \(U\) has weight at most \(w_{S_0}(U)\).
If \(U\) is nef, then \(w_C(U)\le w_{S_0}(U)\) for every fixed
component, and the multiplicity of weight \(mw_{S_0}(U)\) in the
virtual character \(\chi(S,mU)\) is
\(\chi(S_0,mU|_{S_0})\), for \(m>0\).
If \(U\) is nef, on a smooth equivariant resolution \(d:\widehat S\to S\) of the
source retraction \(p:\widehat S\to S_0\), there is an equivariant
nonzero bundle map \(p^*(U|_{S_0})\to d^*U\).
\end{lemma}

\begin{proof}
An ample line has a linearized positive power for a projective $\C^*$-action. Indeed the action's map to the Picard variety is constant, since a torus has no nonconstant homomorphism to an abelian variety. The pairs consisting of an element and an isomorphism from the pulled-back line to the line form an algebraic extension of $\C^*$ by the scalar $\C^*$; this extension is a torus and splits. This supplies the needed linearization. The section-weight bound, the comparison of fixed weights, and the index formula are Lemma~\ref{lem:source-character} for $T=\C^*$ and $U$ as given. The source retraction itself is the attracting-cell map of \cite{BB1973}. The additional assertion needed by the harmonic construction is its bundle map, which we now prove.

For the bundle map, identify the generic fibers of \(U\) and its
restriction at the source by transporting a vector along an orbit,
dividing by \(t^{w_{S_0}(U)}\), and taking the limit. The orbit
extension over \(\mathbb A^1\), over the relevant function field,
pulls the line back to an equivariantly trivial line of that weight.
This constructs a multiplicative rational identification, hence the
claimed rational map. If \(U\) is ample, a generated power has enough
source-weight sections to generate its restriction to \(S_0\). Their
pulled-back restrictions map to the original sections by the limit
construction, proving absence of poles. For nef \(U\), a negative
order in the rational map would give a negative order for
\(kU+H\) for large \(k\), contrary to ampleness. This proves regularity
along every prime divisor and hence everywhere by normality. The
trivial action has \(S_0=S\) and the identity map.
\end{proof}

\begin{proposition}[Attainment of the source weight]\label{har:prop:extremal}
Let \(S\) be smooth connected projective rationally connected, with
smoothly semipositive \(L=-K_S\), and fix a \(\mathbb C^*\)-action with
the natural anticanonical linearization. If \(w=w_{S_0}(L)\), then
some \(\ell>0\) has a nonzero section of \(\ell L\) of weight \(\ell w\).
\end{proposition}

\begin{proof}
\smallskip\noindent\textit{Determinant sections with a prescribed weight.}
A point is immediate. Otherwise \(S_0\) is rationally connected,
being a smooth rational image of \(S\). Thus
\(\chi(S_0,\mathcal O_{S_0})=1\). The highest-weight index in
Lemma~\ref{har:lem:source} is a nonzero polynomial in \(m\), so for
unbounded \(m\) a fixed group \(H^q(S,mL)\) contains weight \(mw\).
The smooth coefficient Lefschetz map is equivariant for a
circle-invariant Kähler form. It gives twisted sections
\[
 v_m\in H^0(S,\Omega_S^{\dim S-q}\otimes(m+1)L)
      =H^0(S,\bigwedge^q T_S\otimes mL)
\]
of weight \(mw\). For \(q=0\) they are the desired sections. Otherwise
saturate their generic span in \(\bigwedge^qT_S\), of rank \(e>0\),
and call it \(\mathcal U\). The determinant extraction used earlier
gives nonzero eigensections
\begin{equation}\label{har:eq:detweights}
 d_i\in H^0(S,\det\mathcal U+k_iL),\qquad
 \operatorname{wt}(d_i)=k_iw,\qquad k_i\to\infty.
\end{equation}
If \(D_i\) is the zero divisor, we claim
\[
 eL-c_1(\mathcal U)\text{ pseudoeffective},\qquad
 (e+k_i)L-D_i\text{ pseudoeffective}.
\]
For $q=\dim S$, saturation gives $\mathcal U=\bigwedge^{\dim S}T_S=L$ and $e=1$, so $eL-c_1(\mathcal U)=0$ directly. For $0<q<\dim S$, the cotangent-subsheaf theorem \cite[Theorem~4.1]{LMPTX2023} applies: \(\mathcal U\otimes(-L)\subset\Omega_S^{\dim S-q}\), and
exterior powers are direct summands of tensor powers in characteristic
zero; saturating can only add an effective determinant correction
in the required direction.

\smallskip\noindent\textit{A base on which dominated systems become rigid.}
Call a divisor class dominated if a positive multiple of \(L\) minus
it is pseudoeffective. Maximize the image dimension of linear systems
of dominated effective integral divisors, and choose a full system
\(|P_0|\) realizing this maximum. Rational connectedness gives
\(\operatorname{Pic}^0(S)=0\), so the system may be linearized and
taken equivariantly. Resolve its graph, take Stein factorization,
flatten by the Hilbert-family construction, resolve the base and
normalize the main component. We obtain
\[
 V\longrightarrow W\longrightarrow S,\qquad g:W\to Y,
\]
with smooth projective \(V,Y\), normal equidimensional \(W\), and
connected fibers. Write \(f:V\to Y\), \(\mu:V\to S\),
\(A=\mu^*L\), \(E=K_{V/S}\), and let \(L_W\) be the pullback of \(L\) to \(W\). An ample pullback from \(Y\) is
dominated: the original moving polarization is dominated, and its
pullback to a birational base model is big, hence dominates a small
ample class. Also \(E\) is dominated. Horizontal exceptional primes
have centers in the original base locus and are bounded by pullbacks
of divisors in \(|P_0|\); vertical primes are bounded by sufficiently
ample base pullbacks. These statements prove domination of the whole
finite effective canonical discrepancy.

A line \(B\) on \(V\) is dominated when \(cA-B\) is pseudoeffective
for some \(c>0\). No such line can have two independent sections on
the generic fiber. If it did, tensor by a sufficiently ample base
pullback to globalize them. This new line remains dominated because
the base polarization is dominated. Combine their ratio with the
original base coordinates, using products to put all coordinates
in a single linear system. Its line is still dominated. Push the
effective system and its pseudoeffective domination relation down
to \(S\). The new ratio is nonconstant along the generic fiber, so
the image dimension exceeds the chosen maximum, a contradiction.
In particular generic-fiber sections of both \(jaA\) and
\(E+jaA\) have dimension at most one for every \(a,j>0\), and the
same adjoint assertion for \(j=0\) follows from domination of \(E\).

There are only finitely many dominated horizontal primes on \(V\).
Here \(\operatorname{Pic}^0(V)=0\), since \(V\) is birational to the
rationally connected \(S\). Thus its Picard group has finite rank.
Infinitely many dominated horizontal primes would give, after clearing
torsion, a nontrivial integral linear equivalence among them. Splitting
positive and negative coefficients gives two effective divisors in
the same class and a pencil varying along the generic fiber, contrary
to the preceding paragraph. The pullbacks of \(D_i\) therefore have
horizontal supports in a fixed finite set. Order their multiplicities
along a subsequence. A ratio from \eqref{har:eq:detweights} gives a
rational eigensection \(s\) of \(aL_W\), weight \(aw\), with no
horizontal poles.

\smallskip\noindent\textit{The metrized base line and its source weight.}
Divisorial normalization, as in
Proposition~\ref{har:prop:maximum}, gives
\begin{equation}\label{har:eq:weightbase}
 g^*Q\longrightarrow aL_W,
\end{equation}
with effective zero divisor omitting one lift over every base prime.
Shift the linearization of \(Q\) to make this map equivariant. Since
the canonical exceptional section times \(s^j\) is nonzero generically,
all required adjoint ranks are exactly one. The maximum construction
gives a continuous semipositive metric on \(Q\). Here the base action
may be nontrivial, but full rank one is available and the ordinary
direct-image line theorem applies without taking an invariant summand.

We must still control the character produced by interpolation on the
base. Let \(Y_0\) be the source of its induced action. The original
unbounded eigensections imply
\begin{equation}\label{har:eq:weightinequality}
 w_{Y_0}(Q)\ge aw.
\end{equation}
After fixing \(a\), take an infinite subsequence of \(k_i\) in one
residue class modulo \(a\), still ordered horizontally. Divide by its
first section. The resulting eigensections of \(jaL_W\), for
unbounded \(j\), have weight \(jaw\), no horizontal poles and vertical
poles bounded by one fixed divisor. Clear these poles by an invariant
effective Cartier divisor \(C\) on \(Y\), with its canonical section
of weight zero. Generic rank one and the map
\eqref{har:eq:weightbase} show that each resulting section comes from
\(jQ+C\). It has no residual poles: at a base prime test on a lift
omitted by the zero divisor of \eqref{har:eq:weightbase}. The source
upper bound gives
\(jaw\le j w_{Y_0}(Q)+w_{Y_0}(C)\). Divide by \(j\) and pass to the
unbounded subsequence to obtain \eqref{har:eq:weightinequality}.

\smallskip\noindent\textit{Euler interpolation on the source of the base.}
The continuous semipositive metric makes \(Q\) nef by
Lemma~\ref{har:lem:smoothing}. Resolve the source retraction of \(Y\)
and dominate this resolution by a smooth model \(V'\) of \(V\).
Let \(p:V'\to Y_0\), \(\mu':V'\to S\), \(A'=\mu'^*L\),
\(E'=K_{V'/S}\), and \(Q_0=Q|_{Y_0}\). Lemma~\ref{har:lem:source}
and \eqref{har:eq:weightbase} give
\(p^*Q_0\to aA'\). An ample pullback by \(p\) is dominated by
\(A'\): the pullback of the original ample base polarization is big
on the resolution of \(Y\), and hence dominates an ample class there,
which in turn dominates an ample pullback from \(Y_0\).
Proposition~\ref{har:prop:dominationvanish} applies to every twist
\(jQ_0\), including zero. The action on \(Y_0\) is trivial, and
\(Q_0\) has the constant fiber weight \(w_{Y_0}(Q)\). Let
\(\mathcal B\) be the weight-zero summand of \(p_*\mathcal O(E')\).
Thus sections of \(\mathcal B\otimes Q_0^j\) have precisely weight
\(j w_{Y_0}(Q)\). Vanishing identifies their dimension with the
Euler polynomial for every \(j\ge0\). At zero this polynomial is
positive, because the canonical exceptional section is invariant.
For some \(j>0\) we therefore obtain a nonzero section of
\(E'+jp^*Q_0\) of that weight. Map it into \(E'+jaA'\) and push
down to \(S\). Its weight is at least \(jaw\) by
\eqref{har:eq:weightinequality}, and at most \(jaw\) by the source
bound on \(S\). Equality gives the required extremal section.
\end{proof}

\begin{corollary}[From extremal characters to invariants]\label{har:cor:torus}
Every algebraic torus acting on the rationally connected \(S\) of
Proposition~\ref{har:prop:extremal} has a nonzero invariant section of
some positive power of \(-K_S\).
\end{corollary}

\begin{proof}
For a generic integral one-parameter subgroup \(\lambda\), its source
\(C\) is a torus-fixed component with all normal weights positive
on \(\lambda\). Denote by \(\gamma_C\) the fiber character of \(L\).
Choose a full torus eigenvector in the extremal section space of
Proposition~\ref{har:prop:extremal}. It cannot vanish on \(C\), since
a nonzero leading conormal jet would strictly lower its
\(\lambda\)-weight. Its character is therefore \(\ell\gamma_C\).
The finitely many characters of such sources have convex hull
containing zero. Otherwise a separating functional, approximated by
a generic rational one, would be negative on all of them, whereas
the source for that one-parameter subgroup has
\(\gamma_C(\lambda)>0\), the sum of its positive normal weights.
Use the effective torus quotient if the action has a kernel; a trivial
action is already covered. A rational convex combination of the source
characters is zero. Clear its denominators and the finitely many
section exponents, then multiply the corresponding powers of their
nonzero eigensections. On the integral variety this product is nonzero
and invariant.
\end{proof}
